% =============================================================================
% chapter3.tex — Project Specifications
% Chapter 3 of the Codara PFE Report
% =============================================================================

\chapter{Project Specifications}
\label{ch:specifications}

% ---------------------------------------------------------------------------
\section{Introduction}
\label{sec:ch3-intro}
% ---------------------------------------------------------------------------

A system as technically ambitious as \codara{} cannot be built without a
precise, measurable specification. Vague goals such as ``improve migration
accuracy'' or ``make translation faster'' are not actionable: they provide no
way to verify that the system has met its objectives, no way to prioritise
design trade-offs, and no way to communicate requirements between the academic
supervisor and the industrial host. This chapter translates the problem
statement developed in Chapter~\ref{ch:context} into a structured set of
objectives, functional and non-functional requirements, use cases, and
high-level architecture constraints.

The specification process followed three steps. First, we identified the
stakeholders --- the data engineering team at the host company, the academic
supervisor, and the end users of the deployed web interface --- and elicited
their goals through structured interviews. Second, we derived requirements from
those goals using the MoSCoW (Must/Should/Could/Won't) prioritisation method,
distinguishing requirements that are essential for a viable system from those
that represent quality improvements. Third, we validated the derived requirements
against the theoretical framework from Chapter~\ref{ch:foundations} to confirm
that each requirement is technically achievable within the project's scope and
timeline.

% ---------------------------------------------------------------------------
\section{Project Objectives}
\label{sec:objectives}
% ---------------------------------------------------------------------------

\subsection{Business Objectives}
\label{subsec:business-objectives}

From the perspective of the host company's data engineering team, the primary
business objective is straightforward: reduce the time and cost of migrating
the organisation's SAS codebase to Python. This objective decomposes into three
concrete sub-goals:

\begin{enumerate}
  \item \textbf{Throughput}: the system must process a multi-file SAS project
    (up to ten files, up to ten thousand lines each) end to end within a
    time budget compatible with an interactive engineering workflow. A
    conversion that takes longer than ten minutes provides no practical
    advantage over asking an engineer to rewrite the code manually.

  \item \textbf{Semantic fidelity}: the translated Python code must be
    behaviourally equivalent to the SAS source for the class of constructs
    in scope. ``Plausible but wrong'' translation --- output that looks correct
    but diverges semantically on edge cases --- is worse than no translation at
    all, because it introduces silent bugs that are harder to find than explicit
    syntax errors.

  \item \textbf{Auditability}: every translation must be accompanied by a
    structured report that tells the engineer exactly what was done, how
    confident the system is in each translated block, and which blocks require
    manual review. This report is a non-negotiable requirement for the host
    company's code-review process.
\end{enumerate}

\subsection{Technical Objectives}
\label{subsec:technical-objectives}

The technical objectives are derived from the business objectives and constrained
by the fifteen-week sprint timeline. They are expressed as measurable targets
wherever possible:

\begin{itemize}
  \item \textbf{Boundary detection accuracy}: $\geq 90\%$ on the 721-block
    gold-standard corpus (see Section~\ref{sec:corpus-construction}).
  \item \textbf{Translation success rate}: $\geq 70\%$ of partitions achieving
    a SemantiCheck Score of LIKELY\_CORRECT or better.
  \item \textbf{Syntax-valid merged output}: $\geq 95\%$ of merged Python
    scripts pass syntax validation without manual intervention.
  \item \textbf{RAPTOR retrieval advantage}: $\geq 10\%$ improvement in
    hit-rate@5 over flat-index retrieval on MODERATE/HIGH-complexity partitions.
  \item \textbf{Streaming performance}: a 10,000-line SAS file must be
    processed by the streaming FSM parser in under 2 seconds with peak memory
    below 100 MB.
  \item \textbf{Complexity calibration}: Expected Calibration Error
    $\mathrm{ECE} < 0.08$ on the held-out 20\% evaluation split.
  \item \textbf{Test coverage}: $\geq 300$ unit and integration tests; target
    statement coverage $\geq 80\%$.
  \item \textbf{Z3 formal verification}: at least 4 SMT patterns with an overall
    provability rate $\geq 35\%$ (fraction of translated blocks receiving a formal
    UNSAT proof).
\end{itemize}

\subsection{Scientific Objectives}
\label{subsec:scientific-objectives}

In the Design Science Research framing, \codara{} is not merely a software
product: it is a research artefact that makes testable scientific claims.
The scientific objectives are therefore the five novelty claims that distinguish
this work from prior art:

\begin{enumerate}
  \item Demonstrate that three-tier RAG with RAPTOR hierarchical clustering
    outperforms flat-index RAG for SAS-to-Python migration in a controlled
    ablation study.
  \item Demonstrate that Z3 SMT verification with CEGAR repair catches real
    semantic errors in LLM-generated Python translations that pass both syntax
    and random-input execution testing.
  \item Demonstrate that CDAIS adversarial synthesis surfaces failure modes
    that random fuzzing misses, and provide formal coverage certificates for
    six SAS semantic error classes.
  \item Demonstrate that the four-layer SemantiCheck Score is a better
    predictor of translation correctness than single-signal metrics (CodeBLEU,
    LLM confidence).
  \item Demonstrate that LLM-as-SAS-oracle is a viable approach to behavioural
    testing of SAS translations without a SAS licence.
\end{enumerate}

% ---------------------------------------------------------------------------
\section{Functional Requirements}
\label{sec:functional-requirements}
% ---------------------------------------------------------------------------

[TABLE 3.1 --- Caption: ``Functional requirements classified by MoSCoW
priority.'' --- Columns: ID, Requirement, Priority (Must/Should/Could),
Source (Business/Technical/Scientific).]

The functional requirements are organised around the eight pipeline stages and
the supporting web interface. The Must-have requirements represent the minimum
viable system:

\textbf{File Ingestion and Parsing (FR-01 to FR-06).} The system must accept
one or more SAS source files via a web upload interface (FR-01), detect and
resolve cross-file \texttt{\%INCLUDE} and libref dependencies (FR-02), and
tokenise each file using a streaming FSM parser that handles files up to
10,000 lines without loading the full file into memory (FR-03). It must detect
partition boundaries between DATA steps and PROC steps with at least 80\%
accuracy using deterministic lark/regex rules (FR-04), and fall back to LLM
boundary resolution for ambiguous cases (FR-05). All parsed metadata must be
persisted to the file registry (FR-06).

\textbf{Complexity Assessment and Routing (FR-07 to FR-09).} Each partition
must receive a risk-level classification (LOW/MODERATE/HIGH/UNCERTAIN) from the
calibrated ML complexity scorer (FR-07), a RAPTOR cluster assignment that
enables hierarchical retrieval (FR-08), and a RAG strategy assignment (Static/
GraphRAG/Agentic) from the StrategyAgent (FR-09).

\textbf{Translation (FR-10 to FR-16).} The system must retrieve relevant
SAS--Python examples from the knowledge base using the assigned RAG strategy
(FR-10), apply deterministic translation rules for unambiguous constructs before
invoking the LLM (FR-11), generate a Python translation using the primary LLM
with structured output (FR-12), cross-verify the translation using an
independent LLM call (FR-13), apply Reflexion repair if cross-verification
fails (FR-14), execute the translation in a sandboxed subprocess to catch
runtime errors (FR-15), and escalate to the CEGAR repair loop if Z3
verification finds a counterexample (FR-16).

\textbf{Verification and Scoring (FR-17 to FR-20).} Each translation must
receive a Z3 SMT verification attempt (FR-17), a CDAIS adversarial test run
(FR-18), a SemantiCheck Score combining all four verification layers (FR-19),
and a fault-localised repair prompt if any layer reports failure (FR-20).

\textbf{Merge and Output (FR-21 to FR-23).} Translated partitions must be
merged into a single Python script with consolidated imports (FR-21),
a dependency-ordered execution sequence (FR-22), and a structured HTML/Markdown
conversion report (FR-23).

\textbf{API and User Interface (FR-24 to FR-30).} The system must expose a
documented REST API with JWT authentication (FR-24), a file upload endpoint
(FR-25), a conversion start endpoint (FR-26), a polling endpoint that returns
live stage-level status (FR-27), a download endpoint for the merged script and
report (FR-28), a knowledge-base management interface (FR-29), and an admin
dashboard for audit logs, user management, and system health (FR-30).

% ---------------------------------------------------------------------------
\section{Non-Functional Requirements}
\label{sec:nonfunctional-requirements}
% ---------------------------------------------------------------------------

[TABLE 3.2 --- Caption: ``Non-functional requirements with acceptance criteria.''
--- Columns: ID, Category, Requirement, Acceptance Criterion.]

The non-functional requirements govern the quality attributes of the system
beyond its functional behaviour.

\textbf{Performance.} End-to-end conversion of a single 500-line SAS file must
complete within 120 seconds on a standard cloud VM (NF-01). The streaming parser
must sustain throughput of at least 5,000 lines per second (NF-02). API
endpoints must respond within 200 ms for all non-pipeline requests (NF-03).

\textbf{Reliability.} The pipeline must not crash on malformed SAS input;
all node failures must be caught, logged, and propagated to the frontend as
stage-level warnings rather than unhandled exceptions (NF-04). Redis checkpoint
recovery must resume a failed pipeline from the last checkpoint without full
re-execution (NF-05).

\textbf{Security.} All API endpoints except \texttt{/health} must require a
valid JWT token (NF-06). Passwords must be hashed with bcrypt (NF-07). The
\texttt{/login} and \texttt{/register} endpoints must be rate-limited to five
attempts per IP per 60 seconds (NF-08). The validation sandbox must isolate
\texttt{exec()} calls in a separate \texttt{multiprocessing.Process} with
removed builtins (NF-09).

\textbf{Maintainability.} All Python modules must follow the project's coding
conventions: \texttt{from \_\_future\_\_ import annotations}, structlog for
logging, Pydantic v2 for data models (NF-10). Test coverage must remain at or
above 80\% as measured by pytest-cov (NF-11). All LLM calls must be logged to
DuckDB for cost tracking and audit (NF-12).

\textbf{Deployability.} The system must ship as a Docker Compose configuration
that starts all services (Redis, backend, frontend) with a single
\texttt{docker compose up} command (NF-13). A GitHub Actions pipeline must run
all tests, benchmarks, and a Docker build check on every push to main (NF-14).

% ---------------------------------------------------------------------------
\section{Use Cases and Scenarios}
\label{sec:use-cases}
% ---------------------------------------------------------------------------

\subsection{Actor Identification}
\label{subsec:actors}

The system distinguishes three actor types:

\begin{itemize}
  \item \textbf{Data Engineer}: the primary user. Uploads SAS files, starts
    conversions, monitors progress, downloads the translated output, submits
    human corrections to improve the knowledge base.
  \item \textbf{Administrator}: manages users, views audit logs, monitors
    system health, and configures pipeline parameters. Has access to all
    Data Engineer functions plus the admin dashboard.
  \item \textbf{System (background)}: the LangGraph pipeline itself, triggered
    by the \texttt{BackgroundTasks} mechanism when a conversion is started.
    Executes asynchronously and writes progress to the database, which the
    Data Engineer polls via the frontend.
\end{itemize}

\subsection{Use Case Diagram}
\label{subsec:use-case-diagram}

[FIGURE 3.1 --- Caption: ``Use case diagram for \codara{}: Data Engineer,
Administrator, and System (pipeline) actors with their associated use cases.''
--- Suggested: UML use case diagram with three actors and approximately twelve
use cases.]

\subsection{Detailed Use Cases}
\label{subsec:detailed-use-cases}

\textbf{UC-01: Upload SAS Files.}
\begin{itemize}
  \item \textit{Actor}: Data Engineer.
  \item \textit{Precondition}: User is authenticated (valid JWT).
  \item \textit{Main flow}: User selects one or more \texttt{.sas} files
    through the web interface. The API stores each file under
    \texttt{backend/uploads/file-\{hex8\}/} (or Azure Blob if configured),
    creates a \texttt{SasFile} record, and returns a list of file metadata
    objects with assigned UUIDs.
  \item \textit{Postcondition}: Files are stored and file records exist in
    the database.
\end{itemize}

\textbf{UC-02: Start Conversion.}
\begin{itemize}
  \item \textit{Actor}: Data Engineer.
  \item \textit{Precondition}: At least one file uploaded (UC-01).
  \item \textit{Main flow}: User selects the target runtime
    (Python or PySpark) and clicks Start. The API creates a
    \texttt{ConversionRow} with status \texttt{queued}, enqueues a job via
    \texttt{PipelineQueueService}, and immediately returns the conversion ID.
    The pipeline begins executing asynchronously.
  \item \textit{Alternative}: If Azure Queue is not configured, the pipeline
    is triggered via \texttt{FastAPI BackgroundTasks} instead.
\end{itemize}

\textbf{UC-03: Monitor Progress.}
\begin{itemize}
  \item \textit{Actor}: Data Engineer.
  \item \textit{Main flow}: The frontend polls \texttt{GET /api/conversions/\{id\}}
    every 1,200 ms. The API returns the conversion's current status and a list of
    \texttt{ConversionStageRow} records, one per pipeline stage, each with its
    own status (pending / running / completed / failed / skipped), latency in
    milliseconds, retry count, and any warnings.
  \item \textit{Termination}: Polling stops when the top-level status is
    \texttt{completed}, \texttt{partial}, or \texttt{failed}.
\end{itemize}

\textbf{UC-04: Download Results.}
\begin{itemize}
  \item \textit{Actor}: Data Engineer.
  \item \textit{Precondition}: Conversion status is \texttt{completed} or
    \texttt{partial}.
  \item \textit{Main flow}: User clicks Download. The API assembles a ZIP
    archive containing the merged Python script and the HTML/Markdown conversion
    report, and returns it as a binary stream.
\end{itemize}

\textbf{UC-05: Submit Correction.}
\begin{itemize}
  \item \textit{Actor}: Data Engineer.
  \item \textit{Precondition}: Conversion completed.
  \item \textit{Main flow}: User provides a corrected Python translation and
    an explanation of the error. The API stores the correction and triggers a
    \texttt{FeedbackIngestionAgent} run that writes the new SAS--Python pair
    to the LanceDB knowledge base, increments the KB version counter, and logs
    the change to the \texttt{kb\_changelog} table.
\end{itemize}

\subsection{Conversion Scenarios}
\label{subsec:conversion-scenarios}

To make the requirements concrete, we define three representative conversion
scenarios that will serve as evaluation cases in Chapter~\ref{ch:implementation}:

\textbf{Scenario A --- Low-Complexity, Static RAG.} A 30-line SAS DATA step
that reads a dataset, computes a new variable via simple arithmetic, and writes
the result. No BY-group processing, no macro references, no cross-file
dependencies. Expected path: \texttt{file\_process} → \texttt{streaming} →
\texttt{chunking} → \texttt{raptor} → risk=LOW → Static RAG → deterministic
rule short-circuit → syntax-valid Python in under 15 seconds.

\textbf{Scenario B --- Cross-File, Graph RAG.} A 150-line SAS programme that
\texttt{\%INCLUDEs} a macro library defined in a second file and applies the
macro to three datasets. Expected path: \texttt{file\_process} detects the
cross-file dependency; \texttt{risk\_routing} assigns MODERATE; GraphRAG
traverses the dependency edge to retrieve the macro library context; the LLM
generates Python with an equivalent helper function.

\textbf{Scenario C --- High-Complexity, Agentic RAG + CEGAR.} A 200-line
SAS programme with RETAIN-based BY-group running totals, a correlated PROC SQL
subquery, and a macro \texttt{\%DO} loop. Expected path: risk=HIGH → Agentic
RAG with escalated \emph{k} → LLM translation → sandbox execution → Z3
counterexample found (RETAIN reset bug) → CEGAR repair loop → verified
translation.

% ---------------------------------------------------------------------------
\section{High-Level Architecture}
\label{sec:high-level-arch}
% ---------------------------------------------------------------------------

The high-level architecture of \codara{} follows a three-tier layered model:

\begin{enumerate}
  \item \textbf{Presentation tier}: a React/Vite single-page application
    running on port 5173 in development (or served by nginx on port 8080 in
    production). All API calls are routed through Vite's proxy to the backend
    at port 8000.
  \item \textbf{Application tier}: a FastAPI ASGI application running on port
    8000, backed by a SQLite operational database, a Redis checkpointing store,
    and an Azure Queue/Blob storage layer for cloud-native job dispatch and
    file persistence.
  \item \textbf{Processing tier}: the eight-node LangGraph pipeline that
    executes asynchronously in a background process. The pipeline reads from
    and writes to a LanceDB vector store (768-dimensional Nomic embeddings),
    a DuckDB analytical store (LLM audit logs), and the SQLite API database
    (stage-level status updates polled by the frontend).
\end{enumerate}

[FIGURE 3.2 --- Caption: ``High-level three-tier architecture of \codara{}:
presentation, application, and processing tiers with their respective data
stores.'' --- Suggested: vertical three-tier diagram with the React frontend
at top, FastAPI in the middle, and the LangGraph pipeline + multi-store data
layer at the bottom.]

The separation between the application tier and the processing tier is
deliberate: it allows the REST API to remain responsive to polling requests
while the pipeline executes a potentially long-running conversion in the
background. This design also makes the processing tier independently scalable:
in a future cloud deployment, multiple pipeline workers could consume from the
Azure Queue in parallel, processing different conversions simultaneously without
any change to the API layer.

% ---------------------------------------------------------------------------
\section{Constraints and Assumptions}
\label{sec:constraints}
% ---------------------------------------------------------------------------

\textbf{Constraints.}
\begin{itemize}
  \item The system must run on a single machine during development. The
    production deployment targets Azure Container Apps but must also work
    locally with Docker Compose for offline development.
  \item All LLM costs must be minimised: deterministic rules and static RAG
    must be preferred over LLM calls whenever the translation is unambiguous.
  \item The system must not require a SAS licence: all evaluation, testing,
    and verification must be achievable with Python-only tooling.
  \item Python 3.11+ is required for the backend; the frontend requires
    Node.js 18+ with Bun as the package manager.
\end{itemize}

\textbf{Assumptions.}
\begin{itemize}
  \item Input SAS files are syntactically valid SAS 9.4 programs (or are at
    least parseable by the lark grammar). Corrupted or partial files produce
    graceful degradation but are not the primary use case.
  \item The LLM providers (Azure OpenAI, Ollama, Groq) are available at
    inference time. The fallback chain handles transient outages but cannot
    function with all three providers simultaneously unavailable.
  \item PySpark is treated as an opt-in target; the frontend type system
    currently exposes only Python, and PySpark support is a future extension.
\end{itemize}

% ---------------------------------------------------------------------------
\section{Conclusion}
\label{sec:ch3-conclusion}
% ---------------------------------------------------------------------------

This chapter has translated the qualitative problem statement of
Chapter~\ref{ch:context} into a precise, testable specification. The eight
measurable technical targets --- boundary detection accuracy, translation
success rate, syntax-valid output rate, RAPTOR retrieval advantage, streaming
performance, ECE calibration, test count, and Z3 coverage --- provide the
evaluation framework used in Chapter~\ref{ch:evaluation}. The three conversion
scenarios (A, B, C) define the functional test paths that are exercised in
Chapter~\ref{ch:implementation}. The MoSCoW-prioritised functional and
non-functional requirements served as the acceptance criteria against which each
sprint's deliverables were validated throughout the development process.

With the specification established, the next chapter presents the architectural
and design decisions made to meet these requirements.
